\begin{document}
\docTitle{Related work, citations, and bibliographies}
\phantomsection
\label{doc:related-work}

\agentProposedEdit
  {agent-add-c8867e37-55e5-514b-8ffb-b991477f0c86}
  {Add this agent-authored block for review.}
  {}
  {The historical shared-bibliography passages in this guide are superseded by the self-contained project policy in \TexLink{docs-template/notes_system.tex}{doc:notes-system}{the notes system}.
Every project that uses citations owns \path{tex/support/references.bib}, while an optional machine-produced override belongs at \path{tex/generated/references.bib}.
The files under \path{tools/latex/} are scaffolding templates and are not compilation dependencies of active document projects.}

This guide defines the standard way to describe related work, cite it in the text, and generate a formal References section in a \code{docs-*} project.
The important distinction is that a related-work note, a citation marker, a bibliography entry, and a citation count are four different things.

\section*{The four pieces}

\begin{enumerate}
  \item The \textbf{Related work section} is the standard related work of any paper that uses \code{\citep} to cite there.
  \item A \textbf{Extended related work} is a long per paper description. 
  \item A full \code{paper_name.tex} is a full paper description. 
  In case of replication we use this approach.
  \item A \textbf{BibTeX entry} in \path{tools/latex/tex/support/references.bib} stores the formal bibliographic metadata shared by manual-document projects.
\end{enumerate}

\section*{Required files}

\agentProposedEdit
  {agent-add-41e9e8f7-da79-50b3-802a-12e69d9a3444}
  {Add this agent-authored block for review.}
  {}
  {The current required layout uses \path{tex/source/} for standalone entries, \path{tex/fragments/} for included TeX, \path{tex/assets/} for authored binary inputs, \path{tex/support/} for styles and bibliographies, and \path{tex/generated/} for required derived inputs.
Generator programs live under \path{scripts/}, while fetched data, raw results, temporary intermediates, and unused outputs live under the ignored \path{.cache/generation/} tree.
The tree below is historical and must not be used for new or migrated projects.}
A project with conventional citations has this structure:

\begin{DocCode}
docs/
|-- tools/
|   `-- latex/
|       |-- manual_style.sty
|       `-- tex/support/references.bib       # canonical shared database
`-- docs-my-topic/
    |-- generated/
    |   `-- tex/support/references.bib       # optional derived snapshot
    `-- manual/
        |-- docs_order_manifest.txt
        `-- src/
            `-- my_note.tex
\end{DocCode}

Use \path{tools/latex/tex/support/references.bib} as the source of truth.
The standard renderer exposes this database to BibTeX for every manual-document project.

A self-contained export, for example an Overleaf project, may include
\path{generated/tex/support/references.bib}.
This file is a generated copy of the canonical database, not a second source of truth.
Update the shared database first and then refresh the snapshot.
The \code{\manualBibliography} command automatically uses the generated snapshot when it is present and otherwise uses the shared database.
Standalone publication projects that own their document class and build pipeline may continue to own a project-specific bibliography.


\section*{Registering a work in BibTeX}

Add one entry for each distinct work to \path{tools/latex/tex/support/references.bib}.
For example:

\begin{DocCode}
@article{kleinebuening2026aligning,
  author  = {Kleine Buening, Thomas and H{\"u}botter, Jonas and
             P{\'a}sztor, Barna and Shenfeld, Idan and
             Ramponi, Giorgia and Krause, Andreas},
  title   = {Aligning Language Models from User Interactions},
  journal = {arXiv preprint arXiv:2603.12273},
  year    = {2026},
  doi     = {10.48550/arXiv.2603.12273},
  url     = {https://arxiv.org/abs/2603.12273}
}
\end{DocCode}

The text before the first comma, here \code{kleinebuening2026aligning}, is the citation key.
Use a stable, descriptive key and do not create a second entry merely because the same work is mentioned in another section or source document.
Prefer metadata from the publisher, proceedings, or official preprint record.

\section*{Extended related work citing mechanism}
\begin{DocCode}
\section*{Extended related work}


\relatedWorkMeta
  {\href{https://arxiv.org/pdf/2603.12273}
    {\emph{Aligning Language Models from User Interactions}}
  \citep{kleinebuening2026aligning}}
  {10 (queried July 11, 2026)}
  {February 2026}
  {ICML 2026 Workshop on Continual Adaptation at Scale}
  {\href{https://arxiv.org/abs/2603.12273}{arXiv}}
  {Thomas Kleine Buening, Andreas Krause}

This paper studies how to align language models directly from multi-turn user interactions. 
It is relevant here because ...
\end{DocCode}

\agentProposedEdit{agent-remove-6a8c126e-c81d-5d9e-9714-40d13e67fa86}{inconsistent}{The five arguments of \code{\relatedWorkMeta} are:}{}
\agentProposedEdit
  {agent-add-74019616-599b-53d0-8131-4bcdd325565a}
  {Add this agent-authored block for review.}
  {}
  {The command takes six arguments: first the linked italic paper title, followed by the five metadata arguments below in their existing order.}

\begin{enumerate}
  \item citation count (i.e. how many citations the paper got so far) and query date;
  \item first-online date;
  \item venue;
  \item repository or source links;
  \item the first author in the author list who is affiliated with a university in Zurich, or \code{-} if none.
\end{enumerate}

The query date is needed because citation counts change over time.

\phantomsection
\section*{Information collection method}
\label{app:related-work-information-method}
We used Codex to retrieve citation counts from Semantic Scholar and first-online dates from arXiv with the following prompt:
\begin{DocCode}
  Find the number of citations for each article using Semantic Scholar, and the first time each article appeared online on arXiv, and update the document.
\end{DocCode}
The first-online metadata records the earliest public online date or source date by month and year.
When a citation count or first-online date could not be found in those sources, the value was completed with ChatGPT and marked with an asterisk.

\section*{Citation commands}

Use the same BibTeX key wherever the work is discussed:

\begin{DocCode}
% Parenthetical numeric citation: "... supervision [1]."
Follow-up interactions can provide hindsight supervision
\citep{kleinebuening2026aligning}.

% Textual citation: "Kleine Buening et al. [1] ..."
\citet{kleinebuening2026aligning} learn from multi-turn
user interactions.

% Several references in one marker, compressed when consecutive.
This observation appears in several studies
\citep{firstKey,secondKey,thirdKey}.
\end{DocCode}

Use \code{\citep} for the usual parenthetical marker and \code{\citet} when the authors are part of the sentence.
Do not type reference numbers such as \code{[1]} manually: numbering can change whenever entries are added or reordered.

\section*{Generating the References section}

Add the following immediately before \code{\end{document}} in every standalone source that uses citations:

\begin{DocCode}
\manualBibliography

\end{document}
\end{DocCode}

\code{\manualBibliography} applies the standard \code{plainnat} bibliography style.
In a normal workspace render, BibTeX resolves \code{tex/support/references.bib} from \path{tools/latex/}.
If \path{generated/tex/support/references.bib} exists in the project, the command uses that self-contained snapshot instead.

Only entries cited by that source appear in its generated References section.
An entry that merely exists in the selected database does not appear.
If an uncited entry must be included intentionally, use:

\begin{DocCode}
\nocite{citationKey}
\end{DocCode}

Prefer a natural citation over \code{\nocite} when the work is actually discussed.

\section*{Complete minimal example}

The source file \path{tex/source/my_note.tex} can be as small as:

\begin{DocCode}
\begin{document}
\docTitle{My research note}

\section*{Related work}

\paragraph{
  \emph{Aligning Language Models from User Interactions}
  \citep{kleinebuening2026aligning}.
}
This work uses follow-up interactions as hindsight feedback.

\manualBibliography

\end{document}
\end{DocCode}

The matching entry must exist in \path{tools/latex/tex/support/references.bib}, or in the derived
\path{generated/tex/support/references.bib} used by a self-contained export.
After rendering, the paragraph contains a resolved numeric citation and the end of the PDF contains the generated References section.

\end{document}
